% Introduction section translated from markdown/2-Introduction.md and refined for Problem C
\section{Introduction}
\subsection{Background}
\textit{Dancing with the Stars} (DWTS) is the U.S. adaptation of the international franchise rooted in the British show \textit{Strictly Come Dancing}. The U.S. version has completed 34 seasons. Each season pairs celebrities with professional dancers who perform weekly. A panel of expert judges assigns technical scores (typically on a 1--10 scale), while viewers vote by phone or online to keep their favorite couple. Fans can vote multiple times up to a weekly limit and cannot vote to eliminate a couple. Judges' scores and fan votes are combined to determine the weekly elimination, and the same combined rule ranks finalists.

In the U.S. version, Seasons 1--2 combined judges and fan votes by rank. Season 2 raised concerns when Jerry Rice reached the finals despite low judges' scores, prompting a change to a percentage-based combination in Seasons 3--27. Another controversy occurred in Season 27 when Bobby Bones won with consistently low judges' scores. Beginning in Season 28, the show introduced a judge save: the bottom two were determined by the combined score, and judges selected which couple to eliminate. Around this time the show also returned to rank-based aggregation; the exact season is unknown, and we follow the problem statement's assumption that it was Season 28. These episodes highlight a persistent tension between technical merit and audience preference.

\subsection{Problem Restatement}
\begin{itemize}
    \item \textbf{Task 1.} Infer the fan vote outcomes from judges' scores and weekly eliminations. Provide two additional metrics: (i) a consistency index between inferred fan votes and actual eliminations; (ii) a certainty analysis for the inferred total fan vote, reported by contestant and by week.
    \item \textbf{Task 2.} Using the Task 1 estimates together with judges' scores, eliminations, and champions, conduct three analyses: (i) compare the ranking method and the percentage method season by season, and determine which is more aligned with audience voting; (ii) examine specific controversial cases (e.g., Jerry Rice in Season 2, Billy Ray Cyrus in Season 4, Bristol Palin in Season 11, Bobby Bones in Season 27); (iii) evaluate how adding a judge-selected bottom-two elimination changes outcomes and recommend which method to use in future seasons.
    \item \textbf{Task 3.} Build a model based on judges' scores, Task 1 inferences, and historical eliminations to evaluate how professional dancers and celebrity attributes (profession, age, industry, and other features) affect outcomes, and compare whether effects are consistent between judges' scores and fan votes.
    \item \textbf{Task 4.} Propose a fairer scoring and elimination system, and justify it with analysis and evidence.
\end{itemize}

\subsection{Our Work}
Judges' scores are effectively on a 10-point scale (with a few values exceeding 10), whereas fan votes can be in the tens or hundreds. Moreover, the ranking and percentage methods operate on different scales. To create a unified basis for analysis, we normalize judges' scores and fan votes within each season-week, which both mitigates outliers and places the two sources on a comparable scale.

For Task 1, we formulate an inverse problem: infer latent fan votes that reproduce observed eliminations. We construct normalized score functions for the ranking and percentage methods, impose elimination-consistency constraints, and optimize a loss function in the fan-vote direction. We report elimination-consistency and cross-week stability indices to keep the inferred votes within a realistic range, and we quantify uncertainty by measuring how sensitive eliminations are to feasible vote perturbations.

For Task 2, we apply both combination rules to every season, compare elimination sets, and quantify which rule better aligns with inferred fan preferences. We then analyze the historical controversies highlighted in the statement and run a counterfactual with a judge-selected bottom-two elimination. These results support our recommendation on which method to adopt going forward and under what conditions a judge save should be applied.

For Task 3, we adapt a multi-head attention (QKV) mechanism to simulate how judges and viewers weight contestant attributes. We perform ablation studies with static feature encoding and semantically enriched feature encoding to validate the attention mechanism. To counter limited data, we orthogonally initialize the Query matrix and introduce a head-diversity penalty. We then analyze the \(\mathrm{softmax}(QK^T)\) weights to interpret feature contributions. Despite regularization (weight decay, cosine annealing) and reduced head dimensions, some overfitting risk remains due to the small sample size.

For Task 4, our analysis suggests the percentage method better reflects audience preference in most weeks, but audience subjectivity intensifies late in the season. We therefore propose a hybrid rule: retain percentage-based aggregation early, and introduce a judge-selected bottom-three elimination starting at mid-season (e.g., Week 7 in a 14-week season, or more generally the first week after the midpoint) to balance professional assessment with audience sentiment.
\par
Tasks 1--2 rely exclusively on the official dataset \texttt{2026\_MCM\_Problem\_C\_Data.csv}. Task 3 augments this with external data (Wikidata, Wikipedia Pageviews, and IMDb) to construct demographic and popularity features, with all sources documented in the Model Preparation section.
